\chapter{The GAMS implementation}\label{chap:gams}

\manref{cge_gams} writes the calibrated model and one scenario as a GAMS
program, so that anyone with GAMS can solve the model with PATH, CONOPT, KNITRO,
IPOPT or another solver; check the results of \pkg with a second, independent
implementation; and extend the equations. The GAMS model is the one that the
package's own tests use to check the Stata model.

\section{What is written}\label{sec:gams-files}

GAMS receives what a user supplies, not the parameters that Stata calibrated:
the SAM after the layout rewrites of \manref{cge_map}
(section~\ref{sec:layouts}); the roles of its accounts; and the settings given
to \manref{cge_calibrate}: elasticities, income elasticities and Frisch
parameters, VAT credit shares, public debt and interest rates, export-only
commodities and export demand, labour-market settings, days lost to illness,
and factor quantities. The GAMS program calibrates every parameter itself,
solves the base, and then solves the scenario. Its equations are an
independent implementation of the equations of this part, written in GAMS
rather than translated from Mata.

\begin{cgetable}{p{0.27\linewidth}p{0.69\linewidth}}
\toprule
File & Content\\
\midrule\endhead
\bottomrule\endlastfoot
\cmd{cge_run.gms} & driver: includes the model, then the scenario\\
\cmd{cge_model.gms} & data, calibration, equations and base solve\\
\cmd{cge_closure.gms} & bounds, variables that do not exist, and the closure\\
\cmd{cge_solve.gms} & the solve statement, for the chosen solver and model type\\
\cmd{cge_write.gms} & writes a solution\\
\cmd{cge_data.inc} & the SAM, the roles of its accounts and the calibration settings\\
\cmd{cge_shock_par.inc} & parameter changes and the closure\\
\cmd{cge_shock_var.inc} & changes to fixed variables\\
\cmd{cge_constraints.inc} & constraints, as complementarity pairs, and the model \cmd{cgecon}\\
\cmd{cge_user.inc}, \cmd{cge_user_write.inc} & user variables and equations, and their results (chapter~\ref{chap:user})\\
\end{cgetable}

The scenario is either the shocks and closure in memory (what the next
\cmd{cge_solve} would solve) or those of a stored simulation. Changed
parameters and fixed variables are written as ratios to their base values, for
example
\begin{filecode}
tm('com_man') = 0 * tm('com_man');
QFS.fx('lab') = 1.05 * QFS.l('lab');
\end{filecode}
so that GAMS applies them to its own calibrated values. A parameter whose base
value is zero is set to its new value. Numbers are written with the fewest
significant digits (15, 16 or 17) that give back the same double. GAMS writes
\cmd{cge_base.csv} and \cmd{cge_}\argu{name}\cmd{.csv}, with one row per
variable and element (scenario, variable, element1, element2, value) and the
model status in the first row.

The equations of \cmd{cge_model.gms} carry the labels of
chapter~\ref{chap:equations} with a prefix for their block: \cmd{EP} for prices
(\cmd{EP1} is \eqref{eq:P1}), \cmd{EQ} for production, \cmd{ET} for trade,
\cmd{EI} for institutions, \cmd{ES} for the system constraints, \cmd{EH} for the
health links and \cmd{EL} for the labour-market regimes. To add an equation,
declare a new variable and its equation after the existing ones, fix one
variable less or free one more so that the system stays square, and give the
new variable a starting level before the base solve.

\textbf{Dynamic runs.} Every period of a run of \manref{cge_dynamic} can be
exported. Each period's changes from the period before are written as ratios,
to parameters (for example capital accumulation, population growth of
subsistence quantities and of the base of the health link, and growth of
government and foreign flows) and to fixed variables, in files named after
the period's simulation: \cmd{cge_dn1_par.inc} and \cmd{cge_dn1_var.inc} for
the simulation \cmd{dn1}. The driver solves the base
and then the periods in order, each from the solution of the period before.
Runs with constraints, microsimulation and sensitivity analysis are not
exported.

\section{Solvers and model types}\label{sec:gams-types}

\begin{cgetable}{p{0.12\linewidth}p{0.55\linewidth}p{0.27\linewidth}}
\toprule
\cmd{type()} & GAMS model & Solvers\\
\midrule\endhead
\bottomrule\endlastfoot
\cmd{cns} (default) & a constrained nonlinear system: the equations, square, with no objective & PATH, CONOPT, KNITRO\\
\cmd{nlp} & the same equations and a constant objective, \cmd{OBJ =e= 0} & IPOPT, CONOPT, KNITRO, MINOS, SNOPT\\
\cmd{mcp} (the default with constraints) & a mixed complementarity problem: the equations, unpaired, and each constraint's bound paired with a slack & PATH\\
\end{cgetable}

\textbf{Constraints.} A scenario with constraints
(chapter~\ref{chap:constraints}) can only be written as a mixed complementarity
problem. For constraint $k$, \cmd{cge_constraints.inc} declares a slack
$\V{SLK}_k \ge 0$; the bound $\V{BND}_k$, $\sigma_k(b_k - Y_k)/s^Y_k \ge 0$,
paired with the slack; and the link $\V{LNK}_k$,
$X_k = \bar x_k + \delta_k\,s^X_k\,\V{SLK}_k$, an unpaired equality that frees
the adjusting variable. When the slack is zero, $Y_k$ may lie within its bound
and $X_k$ stays at $\bar x_k$; when it is positive, $Y_k$ is at its bound. This
is the condition \eqref{eq:B1}, which GAMS passes to PATH exactly, without
smoothing. The variables with no pair must be free or fixed, so the bounds that
keep prices and quantities positive in the CNS form are left out. The base is
solved as the model \cmd{cgebase}, the scenario as \cmd{cgecon}.

\textbf{Tolerances.} Convergence tolerances in GAMS solvers are absolute, so
they are scaled to the largest SAM cell, within what double precision allows:
$p = \max(10^{-11}, 10^{-13} \max_{rc} |\SAM(r, c)|)$. When it runs, the GAMS
program writes an option file for four solvers: PATH (\cmd{convergence_tolerance} $= p$), CONOPT
(\cmd{rtnwma}, $p$ relative to the largest cell), IPOPT (\cmd{tol} $= 10^{-12}$,
\cmd{constr_viol_tol} $= p$) and KNITRO (\cmd{feastol_abs} $= p$,
\cmd{feastol} $= 10^{-12}$). Other solvers use their defaults.

\textbf{Running GAMS.} In the folder written, \cmd{gams cge_run.gms} solves the
base and the scenario with PATH; \texttt{-{}-cgesolver=conopt} chooses another
CNS solver, and \texttt{-{}-cgesolver=ipopt -{}-cgetype=nlp} an NLP solver. From
interactive Stata, \cmd{cge_gams} can start GAMS itself; Stata for Windows
cannot start other programs in batch mode. A GAMS licence is needed; GAMS runs
solvers outside the licence only for small models. Listing files name the
licence holder.

\section{Comparison and import}\label{sec:gams-import}

\textbf{Comparison.} \cmd{compare} reads the GAMS results and computes, for
every variable and element that exists in the model,
\begin{equation*}
\frac{|x^{\text{GAMS}} - x^{\text{Stata}}|}{\max(1, |x^{\text{Stata}}|)},
\end{equation*}
and reports the largest value, for the base and for the scenario, with the
variable where it occurs.

\textbf{The GAMS solution in Stata.} \cmd{import()} stores the GAMS solution as
a Stata simulation: the free variables take the values of GAMS, the fixed ones
those of Stata, and Stata evaluates the model's equations at the result and
reports the largest scaled residual. For a model too large for Stata's dense
solver (\secref{sec:performance}) the route is to build, calibrate and shock
the model in Stata, export it, solve it with PATH, and import the solution for
reports, microsimulation or comparison.

\textbf{Agreement.} On 31 exported cases, which together cover every layout,
extension and closure of the package, the solutions of GAMS agree with those of
Stata to these relative differences:

\begin{cgetable}{lllll}
\toprule
Solver & Type & Cases solved & Base & Scenario\\
\midrule\endhead
\bottomrule\endlastfoot
PATH & CNS & 26 of 26 & $4.3 \times 10^{-16}$ & $1.6 \times 10^{-11}$\\
PATH & MCP, with constraints & 5 of 5 & $2.8 \times 10^{-17}$ & $4.0 \times 10^{-12}$\\
CONOPT & CNS & 17 of 17 & $4.3 \times 10^{-16}$ & $3.6 \times 10^{-9}$\\
KNITRO & CNS & 17 of 17 & $4.3 \times 10^{-16}$ & $1.6 \times 10^{-11}$\\
IPOPT & NLP & 17 of 17 & $5.6 \times 10^{-15}$ & $1.6 \times 10^{-11}$\\
\end{cgetable}

Stata solved every case with a tolerance of $10^{-12}$ on the scaled
residuals, so differences of about $10^{-11}$ are at the precision of the two
solutions; CONOPT stops at a looser feasibility tolerance. CONOPT, KNITRO and
IPOPT do not solve complementarity problems, so they were not run on the cases
with constraints.
